% Mycelium Report Template — OVERVIEW + SUPPLEMENT shape (DEFAULT)
% =============================================================================
% Use when: main text is the overview a collaborator could read in 10 minutes,
% supplement carries methods detail, worked examples for failure modes, and
% exhaustive tables. This is the default shape — it produces the cleanest
% skim-reader experience and naturally absorbs failure-mode worked examples
% in the supplement.
%
% This template renders main + supplement as a single PDF. The supplement
% begins after \appendix and is visually marked as such. If a project prefers
% two separate PDFs, split this template at the \appendix line into
% main.tex and supplement.tex.
%
% Copy to analysis/[name]/reports/[name]-report.tex and replace %%PLACEHOLDER%%
% markers. Numbers and terms must be sourced from .manifest.json; do not
% introduce values that have not been registered there.

\documentclass[11pt,a4paper]{article}

% --- Packages ---
\usepackage[utf8]{inputenc}
\usepackage[T1]{fontenc}
\usepackage{amsmath,amssymb}
\usepackage{graphicx}
\usepackage{booktabs}              % Professional tables
\usepackage{hyperref}
\usepackage[margin=1in]{geometry}
\usepackage{natbib}                % Citations
\usepackage{xcolor}
\usepackage{caption}
\usepackage{subcaption}            % Sub-figures
\usepackage{float}                 % [H] placement
\usepackage{longtable}             % Multi-page tables
\usepackage{listings}              % Code listings

\hypersetup{
    colorlinks=true,
    linkcolor=blue,
    citecolor=blue,
    urlcolor=blue
}

% --- Generated reportable values ---
% Imports macros for every entry in .manifest.json:numbers[*], plus the
% \SciVal{\Macro}{snapshot} and \SciText{\Macro}{snapshot} wrappers.
% Regenerated by `python skills/core/scripts/render_report_values_tex.py
% <path-to-manifest>` at the start of every Phase-7 recompile.
% Verified for drift by scitexlintr before pdflatex runs.
\IfFileExists{build/report_values.tex}{\input{build/report_values.tex}}{%
  % Wrapper fallbacks if the generated file is missing (e.g., legacy reports).
  \providecommand{\SciVal}[2]{#1}%
  \providecommand{\SciText}[2]{#1}%
}

% --- Metadata ---
\title{%%TITLE%%}
\author{%%AUTHOR%%}
\date{%%DATE%%}

\begin{document}

% =============================================================================
% TITLE PAGE
% =============================================================================
\maketitle

% =============================================================================
% ABSTRACT
% =============================================================================
\begin{abstract}
%%ABSTRACT%%
% Lead with the scientific question (from the planning brief), the comparison
% baseline, and the primary metric value sourced from the manifest. State
% the single biggest caveat. No undefined acronyms, no citations, no figure
% references. 150–250 words.
\end{abstract}

\newpage
\tableofcontents
\newpage

% =============================================================================
% MAIN TEXT — the overview a collaborator could read in 10 minutes
% =============================================================================

% =============================================================================
% PROBLEM STATEMENT
% =============================================================================
\section{Problem Statement}

%%PROBLEM_STATEMENT%%
% Frame the question in plain English. Name the comparison baseline. Define
% every concept a non-specialist would need to follow the main text — defer
% the detailed definitions to the supplement.

% =============================================================================
% METHODS (main text — sketch only; full detail in supplement)
% =============================================================================
\section{Methods}

%%METHODS_OVERVIEW%%
% In the overview+supplement shape, main-text Methods is a paragraph:
% 4–6 sentences a non-specialist could follow. Name the inputs, the
% analytical operation, the comparison baseline, and the primary metric.
% For load-bearing methodological choices, lead with the intuitive
% explanation before the technical term — a sentence of intuition before
% the formal name pays off even in the overview.
% Cross-reference the supplement for definitions, parameter rationale, and
% software versions: e.g., "See Supplement Section S1 for definitions and
% Supplement Section S2 for parameter choices."

% =============================================================================
% RESULTS
% =============================================================================
\section{Results}

%%RESULTS%%
% Each result is one to two paragraphs: state the question, the finding
% (with numeric values sourced from the manifest), and the interpretation.
%
% For each new aggregation analysis type, include one HEALTHY worked example
% in the main text (sparkline-style mini-table preferred over a full
% figure). Failure-mode worked examples go in the supplement (S3).
%
% Subsequent claims using the same analysis type reuse the example by
% reference — they do not need their own worked example.
%
% Figure template (use sparingly — main text typically has 2–4 figures total):
% \begin{figure}[H]
%     \centering
%     \includegraphics[width=0.8\textwidth]{../outputs/figures/FIGURE_NAME.pdf}
%     \caption{[What it shows]. [How to read it]. [Key takeaway].}
%     \label{fig:LABEL}
% \end{figure}

% =============================================================================
% CONCLUSIONS
% =============================================================================
\section{Conclusions}

%%CONCLUSIONS%%
% Map each conclusion back to the problem statement. State caveats explicitly
% at a heading level at least as prominent as the claim. Distinguish what
% the data shows from what it suggests. Do not introduce new results.

% =============================================================================
% NEXT STEPS (optional — delete if not applicable)
% =============================================================================
\section{Next Steps}

%%NEXT_STEPS%%
% Only include if there are clear, actionable follow-up analyses. Delete
% this entire section if there is nothing concrete to say.

% =============================================================================
% PROVENANCE
% =============================================================================
\section{Provenance}

\begin{description}
\item[Analysis directory] \texttt{%%ANALYSIS_DIR%%}
\item[Key scripts] ~
  \begin{itemize}
%%SCRIPT_LIST%%
  \end{itemize}
\item[Git commit] \texttt{%%GIT_COMMIT%%}
\item[Analysis period] %%ANALYSIS_PERIOD%%
\item[Analyst] %%ANALYST%%
\item[Report generated] \today
\item[Manifest] \texttt{.manifest.json}
\item[Compile log] \texttt{.compile-log.md}
\item[Software environment] See \texttt{ENVIRONMENTS\_INSTALLATIONS.md}
\end{description}

% =============================================================================
% SUPPLEMENT — methods detail, failure modes, exhaustive tables
% =============================================================================
\appendix
\renewcommand{\thesection}{S\arabic{section}}
\setcounter{section}{0}

\newpage
\begin{center}
{\huge\bfseries Supplement}
\end{center}
\vspace{1em}

% --- S1: Definitions (full) ---
\section{Definitions}

%%SUPPLEMENT_DEFINITIONS%%
% Define every mathematical symbol, technical term, and domain-specific
% concept used anywhere in the report. Group related definitions together.
% For coined terms that shadow established literature terms, add an
% explicit "this is NOT the standard X" disclaimer at first use, naming
% the specific way it differs.

% --- S2: Technical Detail ---
\section{Technical Detail}

%%SUPPLEMENT_TECHNICAL_DETAIL%%
% Datasets, software versions, parameter rationale, statistical methods,
% computational environment. This is the reproducibility surface — someone
% should be able to re-run the analysis from the main-text Methods plus
% this section.

% --- S3: Failure-mode Worked Examples ---
\section{Failure-mode Worked Examples}

%%SUPPLEMENT_FAILURE_MODES%%
% For each new aggregation analysis type whose main-text worked example
% showed a healthy case, include a contrasting failure-mode example here.
% Same format as the main-text example, but the point is to show where
% the analysis breaks.

% --- S4: Parameter Sensitivity ---
\section{Parameter Sensitivity}

%%SUPPLEMENT_SENSITIVITY%%
% Threshold sweeps, method comparisons, subsample stability, random seed
% stability. Each figure has a caption explaining what was varied and the
% conclusion. Referenced from main-text Methods.

% --- S5: Extended Tables ---
\section{Extended Tables}

%%SUPPLEMENT_TABLES%%
% Long tables that would interrupt the main narrative. Use longtable for
% multi-page tables.

% =============================================================================
% REFERENCES
% =============================================================================
\bibliographystyle{plainnat}
\bibliography{references}

\end{document}
